%% kapa: exam=1 round=25 year=2014 period=1 form=A subject=civil_law subject_ko=민법 seq=23 number=23
%% source: Q-Net 감정평가사 기출문제 (cst00309) artlSeq=5101207 file=extracted/1008/5101207/2014년 제25회 감정평가사 1차시험 기출문제/1교시(A형).hwp
%% conversion: kordoc 4.24.51 -> tools/gen1.py (auto); review=proofread-sonnet; flags=none
\begin{question}{23}
조건과 기한에 관한 설명으로 옳지 않은 것은? (다툼이 있으면 판례에 의함)
\begin{choices}
  \item 법률행위에 정지조건이 붙어 있다는 사실의 증명책임은 그 법률효과의 발생을 다투는 자에게 있다.
  \item 조건은 사적 자치에 의한 것으로 당사자가 그 의사에 의하여 임의로 정한 것이어야 한다.
  \item 조건 성취로 불이익을 받을 당사자가 과실로 신의성실에 반하여 조건 성취를 방해한 경우에도 상대방은 그 조건의 성취를 주장할 수 있다.
  \item 기한의 이익은, 당사자의 특약이나 법률행위의 성질에 의하여 분명하지 아니한 경우, 채권자의 이익을 위한 것으로 추정된다.
  \item 당사자가 조건 성취의 효력을 그 성취 전에 소급하게 할 의사를 표시한 경우에는 그 의사에 의한다.
\end{choices}
\end{question}
\answer{4}
